\فصل{نتایج}
\label{chap:results}

\قسمت{تشکیل آرایش چندذره‌ای با طراحی رهاسازی}
\label{sec:results-particle-formation}

آرایش ذرات در کانال مستطیلی دوبعدی به‌صورت عددی بررسی شد.
هندسه و مدل ذره مطابق بخش~\ref{sec:theory-particle-formation} هستند.
همهٔ مکان‌ها، زمان‌ها و فاصله‌های این بخش بی‌بعد هستند.
در هر حالت، هدف رسیدن هم‌زمان مراکز ذرات به نقاط مشخص بود.
خطای نهایی و کمینهٔ فاصله طبق رابطهٔ~\eqref{eq:validation-formation-errors} محاسبه شدند.

\FloatBarrier
\زیرقسمت{تشکیل یک آرایش ساده}
\label{sec:results-formation-triangle}

مثلث، نمونهٔ پایه برای بررسی آرایش با سه ذره است.
دو ذره از دیوارهٔ پایین و یک ذره از بالا رها شدند.
مکان رهاسازی هر ذره روی دیوارهٔ خودش مستقل انتخاب شد.
دو ذرهٔ پایین تقریباً در زمان $t^*=0.2500$ رها شدند.
ذرهٔ بالا در زمان $t^*=0.2894$ رها شد.
هر سه مسیر در زمان هدف $T^*=0.3462$ ارزیابی شدند.

دو پنل نخست شکل~\ref{fig:results-formation-basic} مسیرها و مراکز نهایی مثلث را نشان می‌دهند.
این نمونه با تنظیمات اجرای پیشین و وزن‌های ثابت دوباره محاسبه شد.
بیشینهٔ خطای نهایی در میدان شبکه برابر $2.67\times10^{-9}$ بود.
در میدان تحلیلی، این خطا به $9.08\times10^{-7}$ رسید.
کمینهٔ فاصلهٔ مراکز در هر دو بازپخش حدود $0.12990$ بود.
این فاصله از آستانهٔ پذیرش $0.026$ بیشتر بود.

\begin{figure}[htbp]
\centering
\includegraphics[width=0.95\textwidth,height=0.75\textheight,keepaspectratio]{figs/chapter4/formation-basic-cases.pdf}
\caption{مسیر رهاسازی و مراکز نهایی مثلث و حروف \lr{S}، \lr{C} و \lr{T}.}
\label{fig:results-formation-basic}
\end{figure}

\FloatBarrier
\زیرقسمت{تشکیل حروف \lr{S}، \lr{C} و \lr{T} با شش ایستگاه ثابت}
\label{sec:results-formation-letters}

در هر دیواره، سه ایستگاه در $x^*=0.55,\,1.10,\,1.65$ قرار داشتند.
برای هر حرف، از هر ایستگاه یک ذره رها شد.
سرعت اولیه و زمان رهاسازی برای هر ذره جداگانه تعیین شدند.
سه ردیف پایانی شکل~\ref{fig:results-formation-basic} مسیرهای حروف را نشان می‌دهند.
مراکز نهایی در نمای نزدیک با نقاط هدف مقایسه شده‌اند.

جدول~\ref{tab:results-formation-letters} خطای نهایی و کمینهٔ فاصله را گزارش می‌کند.
بیشینهٔ خطای مرجع در میان حروف، متعلق به حرف \lr{S} بود.
حرف \lr{T} در هر دو بازپخش کمترین فاصلهٔ مراکز را داشت.
کمینهٔ فاصلهٔ مرجع این حرف حدود $0.02746$ بود.
هر سه برنامه در بازپخش عددی معیارهای پذیرش را برآورده کردند.

\begin{table}[htbp]
\centering
\small
\setlength{\tabcolsep}{5pt}
\caption{خطای نهایی و کمینهٔ فاصلهٔ بی‌بعد برای سه حرف.}
\label{tab:results-formation-letters}
\begin{tabular}{@{}lccccc@{}}
\hline
حرف & $T^*$ & $E_{\max,\mathrm{PINN}}^*$ & $E_{\max,\mathrm{ref}}^*$
& $d_{\min,\mathrm{PINN}}^*$ & $d_{\min,\mathrm{ref}}^*$ \\
\hline
\lr{S} & $2.3875$ & $2.073\times10^{-8}$ & $3.768\times10^{-4}$ & $0.0335802$ & $0.0336059$ \\
\lr{C} & $1.8674$ & $1.121\times10^{-8}$ & $1.948\times10^{-4}$ & $0.0348533$ & $0.0348257$ \\
\lr{T} & $2.2739$ & $1.141\times10^{-8}$ & $2.877\times10^{-4}$ & $0.0274005$ & $0.0274613$ \\
\hline
\end{tabular}
\end{table}

اندیس \lr{PINN} به بازپخش با میدان شبکه اشاره دارد.
اندیس \lr{ref} بازپخش مستقل با میدان تحلیلی پوازوی را مشخص می‌کند.
مقدار $T^*$ زمان هدف همان آرایش است.

\FloatBarrier
\زیرقسمت{آرایش اردک و بازپخش مستقل برنامه‌های رهاسازی}
\label{sec:results-formation-ducks}

چهار برنامهٔ ذخیره‌شده برای آرایش اردک دوباره بازپخش شدند.
حالت‌های ۲۰ و ۴۲ ذره‌ای از شش ایستگاه ثابت استفاده کردند.
در این حالت‌ها، هر ایستگاه می‌توانست چند ذره رها کند.
در حالت‌های ۲۸ و ۶۴ ذره‌ای، مکان رهاسازی هر ذره آزاد بود.
هر برنامه با میدان شبکه و میدان تحلیلی پوازوی بازپخش شد.
برنامهٔ رهاسازی هر حالت در هر دو بازپخش یکسان بود.

شکل~\ref{fig:results-formation-ducks} آرایش‌ها را در زمان هدف هر برنامه نشان می‌دهد.
دیسک‌ها اندازهٔ ذرات را با مقیاس هندسی کانال نشان می‌دهند.
جدول~\ref{tab:results-formation-ducks} خطای نهایی و کمینهٔ فاصله را مقایسه می‌کند.

\begin{figure}[htbp]
\centering
\includegraphics[width=0.95\textwidth]{figs/chapter4/formation-duck-replays.pdf}
\caption{چهار آرایش اردک در زمان هدف برنامه‌های رهاسازی.}
\label{fig:results-formation-ducks}
\end{figure}

\begin{table}[htbp]
\centering
\small
\setlength{\tabcolsep}{5pt}
\caption{مقایسهٔ بی‌بعد بازپخش شبکه و میدان تحلیلی برای آرایش اردک.}
\label{tab:results-formation-ducks}
\begin{tabular}{@{}clcccc@{}}
\hline
$N$ & نوع رهاسازی & $E_{\max,\mathrm{PINN}}^*$ & $E_{\max,\mathrm{ref}}^*$
& $d_{\min,\mathrm{PINN}}^*$ & $d_{\min,\mathrm{ref}}^*$ \\
\hline
$20$ & شش ایستگاه ثابت & $7.243\times10^{-4}$ & $7.616\times10^{-4}$ & $0.0296120$ & $0.0293893$ \\
$28$ & مکان آزاد & $2.939\times10^{-6}$ & $1.452\times10^{-6}$ & $0.0286904$ & $0.0286904$ \\
$42$ & شش ایستگاه ثابت & $6.503\times10^{-4}$ & $6.621\times10^{-4}$ & $0.0262260$ & $0.0260481$ \\
$64$ & مکان آزاد & $3.456\times10^{-4}$ & $3.456\times10^{-4}$ & $0.0262883$ & $0.0262883$ \\
\hline
\end{tabular}
\end{table}

هر چهار برنامه در بازپخش عددی پذیرفته شدند.
آرایش ۴۲ذره‌ای کمترین حاشیهٔ فاصله را در میان چهار برنامه داشت.
در بازپخش مرجع، حاشیهٔ این آرایش حدود $4.81\times10^{-5}$ بود.
این حاشیه از اختلاف کمینهٔ فاصله میان دو بازپخش کوچک‌تر بود.
این اختلاف اهمیت بازپخش مستقل را برای برنامه‌های نزدیک آستانه نشان می‌دهد.

% Keep the final comparison paragraph on the preceding text page.
\enlargethispage{2\baselineskip}
در این چهار نمونه، تعداد ذرات و نوع رهاسازی هم‌زمان تغییر می‌کنند.
بنابراین مقایسهٔ آن‌ها عملکرد چهار برنامهٔ مشخص را توصیف می‌کند.
آرایش‌های گزارش‌شده به زمان هدف هر برنامه تعلق دارند.

\FloatBarrier
\زیرقسمت{اثر تعداد ایستگاه‌های رهاسازی}
\label{sec:results-formation-launcher-count}

یک هدف کپسولی با ۲۴ ذره برای مقایسهٔ ایستگاه‌ها انتخاب شد.
نقاط هدف و بودجهٔ بهینه‌سازی در تمام حالت‌ها یکسان بودند.
تعداد ایستگاه‌های هم‌فاصله در هر دیواره از یک تا شش تغییر کرد.
در حالت آزاد، مکان رهاسازی هر ذره مستقل تعیین شد.
هر پیکربندی با سه شروع بهینه‌سازی بررسی شد.

شکل~\ref{fig:results-formation-launcher-count} خطای نهایی و حاشیهٔ فاصلهٔ هر اجرا را نشان می‌دهد.
حاشیهٔ فاصله، اختلاف کمینهٔ فاصلهٔ مرجع با آستانهٔ $0.026$ است.
با دو ایستگاه، هیچ‌یک از سه برنامه پذیرفته نشد.
با چهار ایستگاه، دو برنامه از سه برنامه پذیرفته شدند.
با شش تا دوازده ایستگاه، هر سه شروع در هر چیدمان پذیرفته شدند.
در حالت آزاد نیز هر سه اجرا پذیرفته شدند.
در این پیمایش، تغییر تعداد ایستگاه‌ها مکان آن‌ها را نیز تغییر داد.

\begin{figure}[htbp]
\centering
\includegraphics[width=0.95\textwidth]{figs/chapter4/formation-launcher-count.pdf}
\caption{خطای نهایی و حاشیهٔ فاصله برای چیدمان‌های مختلف ایستگاه‌ها.}
\label{fig:results-formation-launcher-count}
\end{figure}

\par\medskip
\noindent\textbf{سقف رهاسازی و زمان‌بندی ذرات.}\par\smallskip
اثر سقف رهاسازی با شش ایستگاه و ۲۴ ذره بررسی شد.
هر نقطهٔ نمودار، یک رهاسازی ثبت‌شده را نشان می‌دهد.
تعداد رهاسازی‌های هر ایستگاه از برنامهٔ ذخیره‌شده به دست آمد.
نتیجهٔ پذیرش فقط برای اجراهای تکمیل‌شده گزارش شد.
تعداد ایستگاه‌ها در این مقایسه ثابت بود.
شکل~\ref{fig:results-formation-release-capacity} این مقایسه و برنامهٔ منتخب را نشان می‌دهد.
سقف برنامهٔ منتخب، شش رهاسازی برای هر ایستگاه بود.

\begin{figure}[htbp]
\centering
\includegraphics[width=\textwidth]{figs/chapter4/formation-release-capacity.pdf}
\caption{سقف رهاسازی و زمان‌بندی ذرات در شش ایستگاه ثابت.}
\label{fig:results-formation-release-capacity}
\end{figure}

\FloatBarrier

\FloatBarrier
\زیرقسمت{آرایش‌های صدذره‌ای در ناحیه‌های هدف}
\label{sec:results-formation-dense}

چهار ناحیهٔ هدف شامل کپسول، پیکان، ماهی و کپسول دوحفره‌ای بودند.
برای هر ناحیه، همان صد نقطه در سه پیکربندی رهاسازی استفاده شد.
پیکربندی‌ها شامل مکان آزاد و چیدمان‌های شش و دوازده ایستگاهی بودند.
هر پیکربندی برای هر هدف با یک شروع ارزیابی شد.
هر چهار برنامهٔ رهاسازی آزاد معیارهای بازپخش را برآورده کردند.
هشت اجرای متناظر با ایستگاه‌های ثابت پذیرفته نشدند.
این نتیجه به چیدمان‌ها و بودجهٔ جست‌وجوی همین آزمون مربوط است.

شکل~\ref{fig:results-formation-dense} دیسک‌های ذرات را در زمان هدف هر برنامه نشان می‌دهد.
پوشش، نسبت سطح اشغال‌شده در ناحیهٔ هدف به مساحت آن است.
نتایج بازپخش مرجع در جدول~\ref{tab:results-formation-dense} آمده‌اند.
ماهی بیشترین خطای نهایی و کمترین حاشیهٔ فاصله را داشت.
پوشش میان شکل‌ها به مساحت و هندسهٔ ناحیه نیز وابسته است.

\begin{figure}[htbp]
\centering
\includegraphics[width=0.95\textwidth]{figs/chapter4/formation-dense-targets.pdf}
\caption{چهار آرایش صدذره‌ای با رهاسازی آزاد در زمان هدف.}
\label{fig:results-formation-dense}
\end{figure}

\begin{table}[htbp]
\centering
\small
\caption{خطا، کمینهٔ فاصله و پوشش چهار آرایش صدذره‌ای پذیرفته‌شده.}
\label{tab:results-formation-dense}
\begin{tabular}{@{}lccc@{}}
\hline
شکل & $E_{\max,\mathrm{ref}}^*$ & $d_{\min,\mathrm{ref}}^*$ & پوشش (\%) \\
\hline
کپسول & $1.21\times10^{-4}$ & $0.026330$ & $25.7$ \\
پیکان & $4.95\times10^{-4}$ & $0.026283$ & $46.5$ \\
ماهی & $1.32\times10^{-3}$ & $0.026197$ & $37.5$ \\
کپسول دوحفره‌ای & $1.21\times10^{-4}$ & $0.026536$ & $28.5$ \\
\hline
\end{tabular}
\end{table}

\FloatBarrier
\زیرقسمت{اثر مکان مرکز آرایش}
\label{sec:results-formation-center}

یک هدف ۱۸ذره‌ای با حفظ هندسه به موقعیت‌های مختلف منتقل شد.
با ایستگاه‌های ثابت، ۲۵ برنامه از ۴۹ موقعیت آزموده‌شده پذیرفته شدند.
برای هر موقعیت، یک شروع بهینه‌سازی اجرا شد.
در $x_c^*=0.52$، هیچ‌یک از هفت ارتفاع آزموده‌شده پذیرفته نشد.
در سه موقعیت طولی انتهایی، هر موقعیت شش پاسخ پذیرفته‌شده داشت.

شکل~\ref{fig:results-formation-center} نقشهٔ حالت ثابت و مقایسهٔ دو روش را نشان می‌دهد.
مقایسهٔ دو روش در هفت موقعیت روی خط میانی انجام شد.
در خط میانی، شش برنامهٔ ثابت و هفت برنامهٔ آزاد پذیرفته شدند.
در بالادست‌ترین نقطهٔ خط میانی، خطای مرجع حالت ثابت $0.152$ بود.
خطای مرجع حالت آزاد در همان نقطه حدود $5.35\times10^{-5}$ بود.

\begin{figure}[htbp]
\centering
\includegraphics[width=0.95\textwidth]{figs/chapter4/formation-target-center.pdf}
\caption{اثر مکان مرکز هدف و مقایسهٔ رهاسازی در خط میانی کانال.}
\label{fig:results-formation-center}
\end{figure}

\FloatBarrier
\زیرقسمت{اثر زاویه و نسبت محورهای آرایش بیضوی}
\label{sec:results-formation-ellipse}

آرایش ۱۲ذره‌ای برای زاویه‌ها و نسبت محورهای مختلف بررسی شد.
در این پیمایش، مرکز هدف و مساحت بیضی اسمی ثابت بودند.
شکل~\ref{fig:results-formation-ellipse} کوواریانس نقاط هدف و مراکز حاصل را مقایسه می‌کند.
نقشهٔ شکل، میانهٔ خطای کوواریانس شروع‌های انجام‌شده را نشان می‌دهد.
از ۵۷ اجرا، ۵۵ اجرا معیارهای پذیرش را برآورده کردند.

میانهٔ خطای نسبی کوواریانس در اجراهای پذیرفته‌شده حدود $0.076\%$ بود.
بیشینهٔ خطای زاویه برای موارد پذیرفته‌شدهٔ دارای جهت تعریف‌شده، $0.107^\circ$ بود.
برای هدف دایره‌ای، جهت محور اصلی تعریف نمی‌شود.

\begin{figure}[htbp]
\centering
\includegraphics[width=0.95\textwidth]{figs/chapter4/formation-ellipse-covariance.pdf}
\caption{مقایسهٔ هندسه و کوواریانس مراکز هدف و مراکز حاصل.}
\label{fig:results-formation-ellipse}
\end{figure}

\FloatBarrier
\زیرقسمت{کمینهٔ فاصله در طول حرکت و اثر جریمهٔ فاصله}
\label{sec:results-formation-clearance}

شکل~\ref{fig:results-formation-clearance} تغییر فاصلهٔ کمینه را در دو آرایش متراکم نشان می‌دهد.
نمونه‌ها شامل اردک ۴۲ذره‌ای و ماهی ۱۰۰ذره‌ای هستند.
نشانگر جداگانه، کمینهٔ پالایش‌شده را در کنار منحنی نمونه‌برداری‌شده مشخص می‌کند.
کمینهٔ مرجع برای اردک و ماهی به‌ترتیب $0.026048$ و $0.026197$ بود.

در آزمونی جداگانه، اثر جریمهٔ فاصله برای یک آرایش شش‌ذره‌ای بررسی شد.
دو روش با سه شروع متناظر و آزمون پذیرش یکسان مقایسه شدند.
هر سه برنامهٔ دارای جریمه پذیرفته شدند.
هر سه برنامهٔ بدون جریمه به‌دلیل فاصلهٔ ناکافی رد شدند.
بیشینهٔ خطای نهایی مرجع در هر شش اجرا کمتر از $1.8\times10^{-4}$ بود.
این مقایسه ضرورت بررسی فاصلهٔ مراکز در طول حرکت را نشان می‌دهد.

سه قاب، جابه‌جایی جفت بحرانی را نزدیک کمینهٔ فاصله نشان می‌دهند.
هر سه قاب از بازپخش مستقل همان برنامه استخراج شدند.
خط افقی، آستانهٔ پذیرش فاصلهٔ مراکز را مشخص می‌کند.
جفت بحرانی اردک، ذرات هفتم و سی‌ام هستند.
بازتاب برنامه و حساسیت رهاسازی در بخش~\ref{sec:appendix-formation} آمده‌اند.

\begin{figure}[htbp]
\centering
\includegraphics[width=\textwidth]{figs/chapter4/formation-clearance-penalty.pdf}
\caption{کمینهٔ فاصله، حرکت جفت بحرانی و اثر جریمهٔ فاصله.}
\label{fig:results-formation-clearance}
\end{figure}

\FloatBarrier
\قسمت{نتایج پیش‌بینی مستقیم با شبکهٔ عملگر عمیق}
\label{sec:results-forward-deeponet}

در این بخش، پاسخ دو شبکهٔ عملگر عمیق بررسی می‌شود.
شبکهٔ جریان، میدان را برای پروفایل‌های معلوم ورودی پیش‌بینی می‌کند.
شبکهٔ مغناطیسی، مسیر ذره را برای تحریک‌های معلوم پیش‌بینی می‌کند.
وزن‌های هر شبکه در همهٔ مقایسه‌های آن ثابت ماندند.

\FloatBarrier
\زیرقسمت{پیش‌بینی میدان جریان برای پروفایل‌های ورودی متفاوت}
\label{sec:results-forward-flow}

میدان جریان در کانال مستطیلی بی‌بعد با ابعاد $5\times1$ بررسی شد.
مقیاس‌ها مطابق رابطهٔ~\eqref{eq:forward-flow-scales} هستند.
هفت پروفایل منتخب، شکل‌های متفاوت ورودی کانال را نشان می‌دهند.
وزن شبکه در همهٔ مقایسه‌ها ثابت ماند.
مرجع با روش اجزای محدود مستقل محاسبه شد.
نمونه‌ها پس از محاسبهٔ خطا از دوازده ورودی انتخاب شدند.
مقایسه روی $7081$ نقطهٔ مشترک انجام شد.
شکل~\ref{fig:results-forward-flow-exploratory-gallery} پروفایل‌های منتخب را نشان می‌دهد.

\begin{figure}[htbp]
\centering
\includegraphics[width=\textwidth]{figs/chapter4/forward-flow-exploratory-gallery.pdf}
\caption{پروفایل‌های بی‌بعد ورودی در هفت نمونهٔ منتخب.}
\label{fig:results-forward-flow-exploratory-gallery}
\end{figure}

جدول~\ref{tab:results-forward-flow-exploratory} خطاهای سرعت و فشار را گزارش می‌کند.
نقاط مقایسه نزدیک ورودی متراکم‌تر بودند.
وزن همهٔ نقاط در محاسبهٔ نرم خطا یکسان بود.
فشار بدون جابه‌جایی مرجع مقایسه شد.
در نمونهٔ دوقله‌ای، خطای بردار سرعت $0.946\%$ بود.
در نمونهٔ نه‌دوره‌ای، این خطا به $2.753\%$ رسید.

\begin{table}[htbp]
\centering
\small
\caption{خطاهای نسبی سرعت و فشار در هفت نمونهٔ اکتشافی.}
\label{tab:results-forward-flow-exploratory}
\begin{tabular}{@{}lccc@{}}
\hline
نمونه & $\mathrm{Re}$ & $E_2(\mathbf U)\,(\%)$ & $E_2(p)\,(\%)$ \\
\hline
\lr{skew\_b3} & $30$ & $1.603$ & $2.374$ \\
\lr{twin\_b3\_r20} & $20$ & $0.946$ & $1.833$ \\
\lr{waves3\_r20} & $20$ & $1.801$ & $5.043$ \\
\lr{lobes4\_r20} & $20$ & $2.928$ & $2.206$ \\
\lr{mixed} & $30$ & $2.328$ & $5.160$ \\
\lr{shoulder} & $30$ & $2.237$ & $5.658$ \\
\lr{nine} & $55$ & $2.753$ & $6.839$ \\
\hline
\end{tabular}
\end{table}

شکل~\ref{fig:results-forward-flow-chain-skew} ورودی نامتقارن و میدان‌های متناظر را نشان می‌دهد.
پنل برش‌ها، تغییر پروفایل سرعت را در امتداد کانال نشان می‌دهد.
نقشهٔ خطا از اختلاف برداری سرعت محاسبه شد.
دو نقشهٔ سرعت از مقیاس رنگ مشترک استفاده می‌کنند.
مقیاس رنگ خطا مستقل است.
نقشه‌ها برای خوانایی کشیده شده‌اند؛ دامنه همان $5\times1$ است.

\begin{figure}[htbp]
\centering
\includegraphics[width=\textwidth]{figs/chapter4/forward-flow-chain-skew-b3.pdf}
\caption{مقایسهٔ میدان، خطای برداری و برش‌های سرعت برای ورودی نامتقارن.}
\label{fig:results-forward-flow-chain-skew}
\end{figure}

شکل~\ref{fig:results-forward-flow-chain-twin} ساختار دوقله‌ای و تحول پایین‌دست آن را نشان می‌دهد.

\begin{figure}[htbp]
\centering
\includegraphics[width=\textwidth]{figs/chapter4/forward-flow-chain-twin-b3-r20.pdf}
\caption{مقایسهٔ میدان، خطای برداری و برش‌های سرعت برای ورودی دوقله‌ای.}
\label{fig:results-forward-flow-chain-twin}
\end{figure}

شکل~\ref{fig:results-forward-flow-chain-nine} به ورودی نه‌دوره‌ای مجموعهٔ اکتشافی اختصاص دارد.
این ورودی با نمونهٔ نه‌دوره‌ای آزمون انتخاب وزن متفاوت است.

\begin{figure}[htbp]
\centering
\includegraphics[width=\textwidth]{figs/chapter4/forward-flow-chain-nine.pdf}
\caption{مقایسهٔ میدان، خطای برداری و برش‌های سرعت برای ورودی نه‌دوره‌ای.}
\label{fig:results-forward-flow-chain-nine}
\end{figure}

آزمون انتخاب وزن، سه معیار سرعت، فشار و مومنتوم را برآورده نکرد.
نتایج کامل آن آزمون در بخش~\ref{sec:appendix-checkpoint} آمده است.
چهار زنجیرهٔ اکتشافی دیگر در بخش~\ref{sec:appendix-forward} ارائه شده‌اند.

\FloatBarrier
\FloatBarrier
\زیرقسمت{پیش‌بینی مسیر ذره تحت تحریک مغناطیسی}
\label{sec:results-forward-magnetic}

عملکرد شبکهٔ \lr{Master10} با دوازده تحریک زمانی بررسی شد.
هندسه، لولهٔ مستقیم با جریان حامل پوازوی بود.
نمای طولی لوله به‌صورت یک ناحیهٔ مستطیلی رسم شده است.
مختصات $X$ و $Z$ به‌ترتیب عرضی و محوری هستند.
این مختصات و زمان $\tau$ مطابق رابطهٔ~\eqref{eq:forward-mag-scales} بی‌بعد شدند.
مقیاس طولی و عرضی به‌ترتیب فاصلهٔ مشخصهٔ منبع و شعاع لوله‌اند.
حل مرجع با روش \lr{DOP853} و معادلهٔ بیش‌میرا محاسبه شد.
مرجع از تابع پیوستهٔ تحریک استفاده کرد.
شبکه در هر پنجره، نمونه‌های زمانی همان تحریک را دریافت کرد.
بیشترین طول پنجرهٔ استنتاج در این مجموعه $0.0125$ بود.

سه تحریک، پاسخ ذره را در شرایط متفاوت نشان می‌دهند.
پالس گاوسی فقط منبع راست را فعال می‌کند.
شکل~\ref{fig:results-forward-magnetic-chain-03} زنجیرهٔ این پاسخ را نشان می‌دهد.
نقشهٔ میدان در مختصات عرضی ثابت رهاسازی محاسبه شد.
جریان حامل در نمای طولی لوله نشان داده شده است.

\begin{figure}[htbp]
\centering
\includegraphics[width=\textwidth]{figs/chapter4/forward-magnetic-chain-03.pdf}
\caption{پالس گاوسی، میدان بی‌بعد و پاسخ ذره.}
\label{fig:results-forward-magnetic-chain-03}
\end{figure}

در نمونهٔ دوم، رهاسازی خارج از محور و تحریک متناوب است.
شکل~\ref{fig:results-forward-magnetic-chain-08} پاسخ متناظر را نشان می‌دهد.

\begin{figure}[htbp]
\centering
\includegraphics[width=\textwidth]{figs/chapter4/forward-magnetic-chain-08.pdf}
\caption{پالس‌های متناوب، میدان بی‌بعد و مسیر ذرهٔ خارج از محور.}
\label{fig:results-forward-magnetic-chain-08}
\end{figure}

نمونهٔ سوم پاسخ به دو تحریک با اختلاف فاز را نشان می‌دهد.
شکل~\ref{fig:results-forward-magnetic-chain-12} این زنجیره را ارائه می‌کند.
قاب مسیر بزرگ‌نمایی شده و مرز آن دیوارهٔ لوله نیست.

\begin{figure}[htbp]
\centering
\includegraphics[width=\textwidth]{figs/chapter4/forward-magnetic-chain-12.pdf}
\caption{تحریک با اختلاف فاز، میدان بی‌بعد و پاسخ ذره.}
\label{fig:results-forward-magnetic-chain-12}
\end{figure}

مسیرهای شبکه و مرجع روی $801$ زمان مشترک مقایسه شدند.
بیشینهٔ خطاهای عرضی و محوری مطابق رابطهٔ~\eqref{eq:validation-forward-magnetic-errors} محاسبه شدند.
جدول~\ref{tab:results-forward-magnetic-errors} بیشینهٔ خطاها را برای هر دوازده حالت نشان می‌دهد.
حد پذیرش عرضی $10^{-2}$ و حد محوری $3.75\times10^{-4}$ بود.
هر دوازده حالت از نظر هر دو مؤلفه پذیرفته شدند.
حساسیت مرجع به نصف‌شدن گام بیشینه نیز بررسی شد.
این گام از $1.5\times10^{-3}$ به $7.5\times10^{-4}$ کاهش یافت.
پاسخ بدون میدان نیز با جواب تحلیلی مقایسه شد.

\begin{table}[htbp]
\centering
\small
\setlength{\tabcolsep}{4pt}
\renewcommand{\arraystretch}{0.85}
\caption{شرایط آغاز و بیشینهٔ خطاهای بی‌بعد مسیر برای دوازده تحریک.}
\label{tab:results-forward-magnetic-errors}
\begin{tabular}{@{}llccccc@{}}
\hline
شناسه & تحریک & $X_0$ & $Z_0$ & $\tau_f$ &
$10^5 E_{X,\infty}$ & $10^5 E_{Z,\infty}$ \\
\hline
$01$ & ثابت راست & $0$ & $-2$ & $2.25$ & $5.017$ & $1.532$ \\
$02$ & ثابت چپ & $0$ & $-2$ & $2.25$ & $0.785$ & $1.161$ \\
$03$ & پالس گاوسی & $0$ & $-2$ & $2.25$ & $3.199$ & $0.208$ \\
$04$ & سینوسی & $0$ & $-2$ & $2.25$ & $5.349$ & $0.700$ \\
$05$ & تغییر منبع & $0$ & $-2$ & $2.25$ & $2.618$ & $1.212$ \\
$06$ & فرکانس متغیر & $0$ & $-2$ & $2.25$ & $2.035$ & $0.789$ \\
$07$ & تحریک متناوب خارج از محور & $-0.30$ & $-1$ & $2.40$ & $5.066$ & $3.257$ \\
$08$ & پالس‌های متناوب & $0.35$ & $-1.125$ & $2.10$ & $1.885$ & $1.192$ \\
$09$ & فرکانس متغیر دو منبع & $0.15$ & $-0.625$ & $2.25$ & $4.460$ & $0.608$ \\
$10$ & تحریک چندفرکانسی & $-0.40$ & $-1.875$ & $1.35$ & $2.358$ & $3.243$ \\
$11$ & تغییر منبع در پایین‌دست & $0.25$ & $0.50$ & $1.50$ & $2.286$ & $0.550$ \\
$12$ & تحریک با اختلاف فاز & $-0.25$ & $-1.50$ & $2.40$ & $0.957$ & $1.352$ \\
\hline
\end{tabular}
\end{table}

بزرگ‌ترین خطای عرضی، $5.349\times10^{-5}$، در حالت سینوسی ثبت شد.
بزرگ‌ترین خطای محوری، $3.257\times10^{-5}$، به رهاسازی خارج از محور تعلق داشت.
شکل~\ref{fig:results-forward-magnetic-errors} تغییر خطای دو مؤلفه را در زمان نشان می‌دهد.
نمودار خطا، تحریک با فرکانس متغیر را نیز شامل می‌شود.
شناسه‌های مقایسه $03$، $08$، $09$ و $12$ هستند.
در این نمودارها، خطای مطلق بر مقیاس ثابت هر مختصات تقسیم شد.
دو زنجیرهٔ تکمیلی در بخش~\ref{sec:appendix-forward} ارائه شده‌اند.
به‌این‌ترتیب، عبور از مبدأ باعث قلهٔ مصنوعی در خطا نمی‌شود.

\begin{figure}[htbp]
\centering
\includegraphics[width=0.94\textwidth]{figs/chapter4/forward-magnetic-errors.pdf}
\caption{خطای زمانی بی‌بعد مسیر در چهار تحریک منتخب.}
\label{fig:results-forward-magnetic-errors}
\end{figure}

خطای گزارش‌شده، اثر نمونه‌برداری، تقریب شبکه و اتصال پنجره‌ها را شامل می‌شود.
این مقایسه به معادلهٔ بیش‌میرا و دو منبع ثابت مربوط است.
حرکت براونی، گرانش و گذرای اینرسی در هر دو پاسخ حذف شدند.
دامنهٔ تحریک در حالت‌های بررسی‌شده از قید بخش~\ref{sec:forward-magnetic-operator} پیروی کرد.

\FloatBarrier
\قسمت{نتایج بازسازی معکوس و رهگیری مسیر}
\label{sec:results-inverse-and-tracking}

ابتدا بازسازی جریان از دو مسیر ذره بررسی می‌شود.
سپس دو روش با زمان و بدون زمان مقایسه می‌شوند.
نتایج رهگیری مغناطیسی به مدل نهم تعلق دارند.

\FloatBarrier
\زیرقسمت{بازسازی جریان از دو مسیر زمان‌دار}
\label{sec:results-inverse-flow-timed}

بازسازی در کانال مستطیلی بی‌بعد با ابعاد $5\times1$ انجام شد.
مسیرهای مصنوعی از جریان مرجع اجزای محدود به دست آمدند.
عدد رینولدز مرجع در هر سه نمونه $30$ بود.
سرعت‌های اولیه و اعداد استوکس در هر نمونه معلوم بودند.
به موقعیت‌ها نویزی با انحراف معیار بی‌بعد $2\times10^{-4}$ افزوده شد.
موقعیت اولیهٔ هر ذره بدون اغتشاش باقی ماند.
پروفایل و میدان مرجع فقط برای ارزیابی پس از برازش استفاده شدند.
در هر سه بازسازی، وزن‌های گام $94000$ ثابت ماندند.

مختصات مکان با ارتفاع کانال مقیاس شدند.
مقیاس زمان $H/U_{\mathrm{ref}}$ است.
پروفایل و میدان با مقیاس مشترک بخش~\ref{sec:inverse-flow-timed} مقایسه شدند.
سرعت میانگین تخمینی با $A=\bar U/U_{\mathrm{ref}}$ گزارش می‌شود.

بازپخش مستقل دو مسیر زمان‌دار در نمونهٔ گوشه‌ای ب بررسی شد.
معادلهٔ حرکت از شرایط اولیه و در میدان تخمینی حل شد.
خطای مکانی در زمان‌های مشترک اندازه‌گیری شد.
شکل~\ref{fig:results-inverse-timed-chain-b} زنجیرهٔ کامل بازسازی را نشان می‌دهد.
ریشهٔ میانگین مربعات فاصله، $0.0628\%$ و $0.0533\%$ ارتفاع کانال بود.
این معیار، ریشهٔ میانگین مربعات فاصلهٔ اقلیدسی است.
مقدارش $\sqrt{2}$ برابر معیار مؤلفه‌ای رابطهٔ~\eqref{eq:validation-inverse-timed-path} است.
خطای پروفایل و میدان جداگانه در بخش~\ref{sec:appendix-inverse} آمده است.
نزدیکی دو مسیر، دقت کل میدان یا یکتایی آن را تضمین نمی‌کند.

\begin{figure}[htbp]
\centering
\includegraphics[width=\textwidth]{figs/chapter4/inverse-timed-chain-b.pdf}
\caption{مسیرهای زمان‌دار، میدان تخمینی و خطای بازپخش نمونهٔ گوشه‌ای ب.}
\label{fig:results-inverse-timed-chain-b}
\end{figure}

\FloatBarrier

\par\medskip
\noindent\textbf{اثر تعداد مسیرهای زمان‌دار بر بازسازی پروفایل ورودی.}\par\smallskip

بازسازی ورودی برای ده پروفایل و پنج تعداد ذره بررسی شد.
عدد رینولدز مرجع $30$ و سرعت میانگین مرجع بی‌بعد برابر یک بود.
در هر حالت، نخستین $N_p$ مسیر از شش مسیر مشترک استفاده شدند.
با افزایش تعداد ذرات، مسیرهای قبلی و نویز آن‌ها ثابت ماندند.
هر ذره $81$ موقعیت در بازهٔ بی‌بعد $0\leq t^*\leq1.2$ داشت.
عدد استوکس همهٔ ذرات $0.42$ بود.
انحراف معیار نویز گاوسی هر مختصه $2\times10^{-4}$ بود.
موقعیت اولیهٔ هر ذره دقیق در نظر گرفته شد.
چیدمان رهاسازی در شکل~\ref{fig:inverse-count-design} پیوست نشان داده شده است.



مسیرهای مصنوعی از میدان اجزای محدود یا حل تحلیلی پوازوی تولید شدند.
مرجع سهموی از حل تحلیلی به دست آمد.
وزن‌های گام $94000$ و تنظیمات منظم‌سازی در همهٔ حالت‌ها ثابت ماندند.
پروفایل ورودی با $65$ گره نمایش داده شد.
هر برازش با دو شروع و سقف $200$ تکرار انجام شد.
پاسخ منتخب با کمترین هزینهٔ برازش میان نامزدهای مجاز تعیین شد.
پروفایل حقیقی برای ارزیابی پس از برازش به کار رفت.
معیار شکل~\ref{fig:inverse-count-errors} خطای نسبی پروفایل در مقیاس مشترک $U_{\mathrm{ref}}$ است.
ترتیب پروفایل‌ها برحسب ریشهٔ میانگین مربعات مشتق دوم ورودی است.

\begin{figure}[htbp]
\centering
\includegraphics[width=0.98\textwidth]{figs/chapter4/inverse-particle-count-errors.pdf}
\caption{خطای بازسازی ورودی برای ده پروفایل و تعدادهای مختلف ذره.}
\label{fig:inverse-count-errors}
\end{figure}

میانهٔ خطا از $21.99\%$ برای دو ذره به $5.31\%$ رسید.
خطای شش‌ذره‌ای در هر ده پروفایل از حالت دوذره‌ای کمتر بود.
از $40$ افزایش تک‌ذره‌ای، $29$ مورد با کاهش خطا همراه بودند.
فقط دو پروفایل در همهٔ افزایش‌ها خطای کمتری داشتند.
در ورودی بتای باریک، خطا ابتدا کاهش یافت و سپس افزایش یافت.
خطای ورودی نه‌نوسانی برای همهٔ تعدادها نزدیک $39\%$ ماند.
خطای ورودی دوقله‌ای با پنج مسیر $1.79\%$ بود.
شکل~\ref{fig:inverse-count-profiles} سه ورودی و سه تعداد مسیر را مقایسه می‌کند.

\begin{figure}[htbp]
\centering
\includegraphics[width=0.95\textwidth]{figs/chapter4/inverse-particle-count-profiles.pdf}
\caption{پروفایل مرجع و بازسازی‌شده برای سه ورودی و سه تعداد مسیر.}
\label{fig:inverse-count-profiles}
\end{figure}

خطای بسیار کوچک سهموی با شکل پایهٔ مقداردهی اولیه نیز مرتبط است.
در این آزمایش، تعداد مشاهدات و پوشش مکانی هم‌زمان تغییر کردند.
این مقایسه به همین ترتیب رهاسازی و تحقق نویز محدود است.
هر $200$ بازپخش به پایان بازهٔ مشاهده رسیدند.
همگرایی بهینه‌ساز در هیچ‌یک از $50$ پاسخ منتخب تأیید نشد.
بنابراین خطاها عملکرد روش را با بودجهٔ تعیین‌شده نشان می‌دهند.

\FloatBarrier
\FloatBarrier
\زیرقسمت{بازسازی جریان از هندسهٔ مسیرها بدون زمان اندازه‌گیری‌شده}
\label{sec:results-inverse-flow-geometry}

همان سه نمونه با استفاده از منحنی‌های مرتب‌شده بررسی شدند.
نقاط منحنی با فاصلهٔ یکنواخت در طول قوس بازنمونه‌برداری شدند.
سرعت اولیهٔ هر ذره همچنان معلوم بود.
زمان‌های مشاهده و پاسخ زمان‌دار وارد برازش هندسی نشدند.
مدت تخمینی، افق بازپخش مستقل هر ذره را تعیین کرد.
هر شش بازپخش در افق تخمینی کامل شدند.

شکل~\ref{fig:results-inverse-geometry-chain-oblique} زنجیرهٔ نمونهٔ مایل را نشان می‌دهد.
جدول~\ref{tab:results-inverse-geometry} پیوست، معیارهای مشترک هر ذره را گزارش می‌کند.
تعریف این معیارها در رابطهٔ~\eqref{eq:inverse-flow-geometry-errors} آمده است.
هر دو بازپخش با یک منحنی مشاهده‌شده مقایسه شدند.
در معیار مرتب، نقاط متناظر در طول قوس نرمال‌شده مقایسه شدند.


\begin{figure}[htbp]
\centering
\includegraphics[width=\textwidth]{figs/chapter4/inverse-geometry-chain-oblique.pdf}
\caption{هندسهٔ مسیر، بازپخش مستقل و زمان‌بندی تخمینی نمونهٔ مایل.}
\label{fig:results-inverse-geometry-chain-oblique}
\end{figure}




در نمونهٔ مایل، هر دو معیار هندسی برای هر ذره کاهش یافتند.
این کاهش در چهار مسیر از شش مسیر مشاهده شد.
خطای متقارن بین $0.0213\%$ و $0.0952\%$ ارتفاع کانال بود.
هر بازپخش برحسب زمان همان روش رسم شد.
در هر دو نمونهٔ گوشه‌ای، خطاهای ذرهٔ دوم افزایش یافتند.
بیشینهٔ فاصلهٔ هاوسدورف در روش بدون زمان حدود $2.62\times10^{-3}$ بود.
اختلاف نسبی پروفایل‌ها در سه نمونه از $9.661\%$ تا $25.265\%$ بود.
مبنای این اختلاف، پروفایل تخمینی روش زمان‌دار است.
با نرمال‌سازی هر پروفایل بر میانگین خودش، اختلاف شکل همچنان دیده شد.
بزرگ‌ترین اختلاف شکل در نمونهٔ مایل، برابر $25.258\%$ بود.

شکل~\ref{fig:results-inverse-duration} مدت تخمینی را با مدت مشاهده مقایسه می‌کند.
در روش زمان‌دار، مدت حرکت از داده معلوم بود.
در روش بدون زمان، این مدت از زمان‌بندی تخمینی به دست آمد.
مدت مشاهده فقط در ارزیابی پس از برازش استفاده شد.
بزرگ‌ترین اختلاف نسبی مدت، $4.393\%$، برای ذرهٔ اول نمونهٔ مایل بود.
در آن مسیر، مدت مشاهده $0.55$ و مدت تخمینی $0.574160$ بود.
این مدت به انتهای منحنی برازش‌شده مربوط است.

\begin{figure}[htbp]
\centering
\includegraphics[width=0.95\textwidth]{figs/chapter4/inverse-geometry-duration.pdf}
\caption{مقایسهٔ مدت بی‌بعد مشاهده و مدت تخمینی هر مسیر.}
\label{fig:results-inverse-duration}
\end{figure}

در روش بدون زمان، هر دو شروع به سقف $700$ تکرار رسیدند.
تنظیمات و بودجهٔ بهینه‌سازی دو روش متفاوت بودند.
بنابراین اختلاف نتایج را نمی‌توان فقط به حذف زمان نسبت داد.

\FloatBarrier
\زیرقسمت{رهگیری مسیر هدف با تحریک مغناطیسی در مدل نهم}
\label{sec:results-master9-tracking}

چهار سناریوی رهگیری در جریان پوازوی لولهٔ مستقیم بررسی شدند.
نمای طولی لوله به‌شکل یک ناحیهٔ مستطیلی نمایش داده می‌شود.
مقیاس‌های $X$، $Z$ و $\tau$ مطابق رابطهٔ~\eqref{eq:forward-mag-scales} هستند.
مسیر هدف با مختصات محوری $Z$ تعریف شد.
گام بی‌بعد انتخاب تحریک برابر $0.025$ بود.
تحریک مؤثر هر گام طبق رابطهٔ~\eqref{eq:inverse-master9-tracking} انتخاب شد.
همان دنبالهٔ تحریک با حل‌گر اینرسی‌دار \lr{Radau} بازپخش شد.

مقیاس طولی و عرضی بر طول‌های مرجع متفاوت استوارند.
بازهٔ نمایش مسیر با دامنهٔ حرکت هر سناریو تنظیم شده است.
فرمان مثبت منبع راست و فرمان منفی منبع چپ را فعال می‌کند.
در فرمان صفر، میدان مغناطیسی خاموش است.
این فرمان، جریان اندازه‌گیری‌شدهٔ سیم‌پیچ نیست.

شکل~\ref{fig:results-master9-tracking} مسیر هدف و بازپخش دو سناریو را نشان می‌دهد.
پنل‌ها تحریک پله‌ای و خطای عرضی رهگیری را نیز نشان می‌دهند.
برای سنجش خطا، هدف در مکان محوری بازپخش ارزیابی شد.
اختلاف دو حل در زمان‌های یکسان محاسبه شد.

\begin{figure}[htbp]
\centering
\includegraphics[width=0.98\textwidth,height=0.72\textheight,keepaspectratio]{figs/chapter4/inverse-master9-tracking.pdf}
\caption{مسیر، تحریک و خطای رهگیری در حالت اسمی و حالت کم‌سرعت.}
\label{fig:results-master9-tracking}
\end{figure}

جدول~\ref{tab:results-master9-tracking} خطا و موقعیت محوری نهایی را گزارش می‌کند.
خطای رهگیری مطابق رابطهٔ~\eqref{eq:validation-master9-tracking} تعریف شد.
آستانهٔ پایان برای مسیر شبکه $Z=1.4$ بود.
شرط تکمیل برای بازپخش، رسیدن به دست‌کم $Z=1.3$ بود.
بازپخش هر چهار سناریو از $Z=1.4$ نیز عبور کرد.

\begin{table}[htbp]
\centering
\small
\setlength{\tabcolsep}{4pt}
\renewcommand{\arraystretch}{0.85}
\caption{خطاهای بی‌بعد و وضعیت تکمیل رهگیری در چهار سناریو.}
\label{tab:results-master9-tracking}
\begin{tabular}{@{}lcccccc@{}}
\hline
سناریو & $\mathrm{Re}$ & $10^4 E_{\mathrm{track},\infty}$ &
$10^4\Delta X_{\max}$ & $10^4\Delta Z_{\max}$ &
$Z_{\mathrm{Radau},f}$ & تکمیل \\
\hline
اسمی & $0.687$ & $0.572$ & $0.595$ & $8.052$ & $1.42266$ & بله \\
سرعت کم و گرانروی زیاد & $0.125$ & $11.529$ & $11.506$ & $2.592$ & $1.40304$ & بله \\
رینولدز زیاد و رهاسازی سریع & $4.950$ & $1.123$ & $1.122$ & $3.009$ & $1.40466$ & بله \\
رهاسازی جابه‌جاشده & $0.7125$ & $0.413$ & $0.451$ & $2.502$ & $1.40170$ & بله \\
\hline
\end{tabular}
\end{table}

$\Delta X_{\max}$ و $\Delta Z_{\max}$ بیشینهٔ اختلاف مکانی دو حل هستند.
حالت کم‌سرعت بیشترین خطای عرضی رهگیری را داشت.
بیشترین اختلاف محوری دو حل در حالت اسمی ثبت شد.
هر چهار سناریو معیارهای عددی رهگیری را برآورده کردند.
این مقایسه برای خانوادهٔ مسیرهای هدف نرم انجام شد.

\FloatBarrier
\FloatBarrier
\زیرقسمت{بازیابی تحریک مغناطیسی از مسیر زمان‌دار ذره}
\label{sec:results-master10-inverse-pending}

موقعیت‌های زمان‌دار مصنوعی برای بازیابی دامنهٔ منابع استفاده شدند.
جریان حامل، هندسه و پارامترهای سیال معلوم و ثابت بودند.
وزن‌های آموزش‌دیدهٔ مدل دهم در تمام برازش‌ها ثابت ماندند.
روش از هدف رابطهٔ~\eqref{eq:inverse-master10-proposed-objective} استفاده می‌کند.
در آزمون تک‌منبعی، تحریک چپ مثبت و تحریک راست صفر بود.
خطای نسبی تحریک تخمینی در این آزمون $0.374\%$ بود.
شکل~\ref{fig:results-master10-single} تحریک و بازپخش مستقل را نشان می‌دهد.
بیشینهٔ خطاهای عرضی و محوری، $6.315\times10^{-5}$ و $1.144\times10^{-5}$ بودند.
بازپخش مستقل با حل‌گر \lr{DOP853} و ضرایب خام انجام شد.
مقیاس‌های عرضی و محوری، طول‌های مرجع متفاوت دارند.


\begin{figure}[htbp]
\centering
\includegraphics[width=\textwidth]{figs/chapter4/inverse-master10-single.pdf}
\caption{بازیابی تحریک منبع چپ و بازپخش مستقل مسیر ذره.}
\label{fig:results-master10-single}
\end{figure}

برازش منتخب تک‌منبعی با توقف جست‌وجوی خطی پایان یافت.
دو شروع دیگر تک‌منبعی همگرا شدند.
در آزمون دومنبعی، نزدیکی مسیر برای بازیابی دو تحریک کافی نبود.
خطای تحریک دومنبعی، پس از هم‌ترازی علامت، $104.65\%$ بود.
شروع منتخب آن به سقف تکرار رسید.
جزئیات حساسیت و نویز در بخش~\ref{sec:appendix-magnetic-inverse} آمده‌اند.
این نتایج، بازیابی یکتای دو تحریک را تأیید نمی‌کنند.

شکل~\ref{fig:results-master10-ambiguity} ابهام علامت پس از فاصلهٔ مشترک صفر را نشان می‌دهد.
دو تحریک در پالس دوم، هم‌زمان علامت عوض می‌کنند.
هیچ علامت کلیِ واحدی دو زوج تحریک را یکسان نمی‌کند.
نیروی فیزیکی و مسیرهای مستقل، یکسان باقی می‌مانند.
اختلاف نسبی تحریک‌ها پس از هم‌ترازی کلی برابر $135.69\%$ بود.
بیشینهٔ اختلاف دو مسیر بی‌بعد برابر $0.00$ بود.
نصف‌کردن گام بیشینهٔ مرجع نیز برای بررسی حساسیت انجام شد.
ترجیح علامت شبکه، شاهد تعیین‌پذیری فیزیکی نیست.
از یک مسیر نزدیک، یکتایی تحریک یا میدان نتیجه نمی‌شود.

\begin{figure}[htbp]
\centering
\includegraphics[width=0.98\textwidth]{figs/chapter4/inverse-master10-ambiguity.pdf}
\caption{تحریک‌های متفاوت پس از فاصلهٔ صفر و مسیرهای مستقل هم‌پوشان.}
\label{fig:results-master10-ambiguity}
\end{figure}
\FloatBarrier
